\section{Error Analysis}
\label{sec:error-analysis}

\subsection{Standing Assumptions}
\label{sec:standing-assumptions}

The estimates in this section are conditional on the following standing assumptions, which we state explicitly rather than leaving implicit.

\begin{itemize}
\item[(A1)] \textbf{Well-posedness.} The optimality system \eqref{eq:state_opt}--\eqref{eq:terminal} admits a unique solution $(y^*,u^*,\lambda^*)$ in an appropriate energy class, and the linearization of this system about $(y^*,\lambda^*)$ satisfies an inf-sup (stability) condition with constant $C_{\text{stab}}$.
\item[(A2)] \textbf{Regularity of the Dirichlet control problem.} Dirichlet boundary control of a parabolic equation is, in general, a low-regularity problem: for boundary data $u\in L^2(\Sigma)$ the state $y$ is typically only a \emph{very weak solution} (transposition solution), and the trace $\partial_n\lambda|_\Sigma$ appearing in \eqref{eq:stationarity} need not exist in the classical sense~\cite{lasiecka2000control,troltzsch2010optimal}. We assume throughout that the data $(y_0, y_d, u^*)$ are smooth enough, and $\Omega$ is regular enough, that $(y^*,\lambda^*)$ lie in $H^{2,1}(Q)$-type spaces in which all traces below are classically well defined. This is a genuine restriction, not a technicality; results that require only $u \in L^2(\Sigma)$ are outside the scope of the estimates below.
\item[(A3)] \textbf{Nonlinearity.} $f\in C^1(\mathbb{R})$ and $f'$ is bounded below, $f'(s)\geq -c_f$ for some $c_f\geq 0$ (satisfied by the Allen--Cahn nonlinearity $f(y)=y^3-y$, for which $f'(s) = 3s^2-1 \geq -1$).
\item[(A4)] \textbf{Approximation capacity.} The network architectures for $Y_\theta, U_\psi, \Lambda_\phi$ are rich enough that the best-approximation error $\epsilon_{\text{net}}$ defined in \eqref{eq:net_approx} can be made arbitrarily small by increasing width/depth. No explicit rate in network size is claimed here; this is treated as an assumption, not a proven property of the architecture.
\end{itemize}

\subsection{Decomposition of the Total Error}
\label{sec:error-decomposition}

Let $e_y := y^* - Y_\theta$, $e_\lambda := \lambda^* - \Lambda_\phi$, $e_u := u^* - U_\psi$ denote the errors in the state, adjoint, and control, relative to the exact solution of the optimality system (not to a discretized/mesh solution). The total error is measured in the norm
\begin{equation}
\| (e_y, e_\lambda, e_u) \|_{\mathcal{E}}^2 := \| e_y \|_{L^2(0,T;H^1(\Omega))}^2 + \| e_\lambda \|_{L^2(0,T;H^1(\Omega))}^2 + \| e_u \|_{L^2(\Sigma)}^2.
\end{equation}
Using the triangle inequality and the stability assumption (A1), the total error can be decomposed into three main contributions:
\begin{equation}
\label{eq:decomp}
\| (e_y, e_\lambda, e_u) \|_{\mathcal{E}} \lesssim \mathcal{E}_{\text{approx}} + \mathcal{E}_{\text{opt}} + \mathcal{E}_{\text{quad}},
\end{equation}
where $\mathcal{E}_{\text{approx}}$ is the best-approximation error of the neural network architecture (item (A4)), $\mathcal{E}_{\text{opt}}$ is the optimization error, i.e.\ how far the trained loss \eqref{eq:indirect_loss} is from its infimum over the chosen architecture, and $\mathcal{E}_{\text{quad}}$ is the quadrature (sampling) error between the Monte-Carlo loss $\mathcal{L}_{\text{ind}}$ and the continuous residual norms it estimates. We do not list a separate boundary or terminal contribution because these conditions are hard-constrained; the state initial mismatch is retained explicitly in the trainable loss and in the estimator of Section~\ref{sec:aposteriori}.

\subsection{A Priori Error Estimate}
\label{sec:apriori}

\begin{theorem}[Conditional a priori estimate]
\label{thm:apriori}
Under Assumptions (A1)--(A3), let $(y^*, u^*, \lambda^*)$ be the exact solution of \eqref{eq:state_opt}--\eqref{eq:terminal}. Suppose the networks $Y_\theta$, $U_\psi$, $\Lambda_\phi$ can approximate $(y^*,\lambda^*,u^*)$ to accuracy $\epsilon_{\text{net}}$ in the $\mathcal{E}$-norm,
\begin{equation}
\label{eq:net_approx}
\inf_{\theta,\psi,\phi} \| (y^* - Y_\theta, \lambda^* - \Lambda_\phi, u^* - U_\psi) \|_{\mathcal{E}} \leq \epsilon_{\text{net}},
\end{equation}
and suppose training reaches parameters $(\theta^*,\psi^*,\phi^*)$ with continuous residual norm
\begin{equation}
\|\mathcal{R}_y\|_{L^2_{\rm av}(Q)}^2 + \|\mathcal{R}_\lambda\|_{L^2_{\rm av}(Q)}^2 + \|\mathcal{R}_u\|_{L^2_{\rm av}(\Sigma)}^2 \;\leq\; \epsilon_{\text{loss}}^2.
\end{equation}
Then
\begin{equation}
\label{eq:apriori_estimate}
\| (y^* - Y_\theta^*, \lambda^* - \Lambda_\phi^*, u^* - U_\psi^*) \|_{\mathcal{E}} \;\lesssim\; \epsilon_{\text{net}} + C_{\text{stab}}\,\epsilon_{\text{loss}},
\end{equation}
with $C_{\text{stab}}$ the stability constant of Assumption (A1).
\end{theorem}

\begin{proof}[Sketch of proof]
Write $e=(e_y,e_\lambda,e_u)$ for the error at the trained parameters, and split $e = e^{\mathrm{approx}} + e^{\mathrm{res}}$, where $e^{\mathrm{approx}}$ is the error of the \emph{best} network approximation (controlled by $\epsilon_{\text{net}}$ via (A4)) and $e^{\mathrm{res}} := (y^*-Y_\theta^*,\dots) - e^{\mathrm{approx}}$ is what remains. By linearity, $e^{\mathrm{res}}$ solves the linearization of \eqref{eq:state_opt}--\eqref{eq:stationarity} about $(y^*,\lambda^*)$ with right-hand side equal to $(\mathcal{R}_y,\mathcal{R}_\lambda,\mathcal{R}_u)$ evaluated at the trained networks (this uses (A2) to make the linearization and its traces meaningful, and (A3) to control the sign of the reaction term $f'(y^*)$ so that the linearized adjoint operator is coercive up to a lower-order term, via a standard exponential-weight / G\r{a}rding-inequality argument). The stability assumption (A1) and the equivalence in \eqref{eq:averaged_residual_norms} then give $\|e^{\mathrm{res}}\|_{\mathcal{E}} \leq C_{\text{stab}}\big(\|\mathcal{R}_y\|_{L^2_{\rm av}(Q)} + \|\mathcal{R}_\lambda\|_{L^2_{\rm av}(Q)} + \|\mathcal{R}_u\|_{L^2_{\rm av}(\Sigma)}\big) \leq C_{\text{stab}}\,\epsilon_{\text{loss}}$, after absorbing the fixed domain factors into $C_{\text{stab}}$. Combining with $\|e^{\mathrm{approx}}\|_{\mathcal{E}}\leq \epsilon_{\text{net}}$ and the triangle inequality gives \eqref{eq:apriori_estimate}. We emphasize that this is a \emph{sketch}: a full proof requires making (A1)--(A2) precise (function spaces, trace theorems) for the specific $f$ and domain considered, which is not carried out here.
\end{proof}

\begin{remark}
\label{rem:continuous-vs-discrete}
Bound \eqref{eq:apriori_estimate} uses the \emph{continuous} residual norm $\epsilon_{\text{loss}}$, not the Monte-Carlo loss $\mathcal{L}_{\text{ind}}$ that is actually minimized during training. The gap between the two is exactly $\mathcal{E}_{\text{quad}}$ in \eqref{eq:decomp}; see Section~\ref{sec:practical} for how it is controlled in practice.
\end{remark}

\begin{remark}
For the Allen--Cahn nonlinearity $f(y)=y^3-y$, $f'(s)=3s^2-1\geq-1$ is bounded below, which is what makes the coercivity argument in the proof sketch available; a genuine coercivity estimate for the linearized adjoint operator additionally requires either a smallness condition on $T$ or an exponential time-weight $e^{-c_f t}$ to absorb the $-1$ reaction term, which we do not carry out in detail here.
\end{remark}

\subsection{Computable A Posteriori Bound}
\label{sec:aposteriori}

Because $Y_\theta, U_\psi, \Lambda_\phi$ are smooth functions on all of $\overline{Q}$ (for smooth activation functions), there is no computational mesh and no notion of an inter-element jump in the normal derivative -- such jumps are identically zero for a globally differentiable network. We therefore build the a posteriori indicator directly from pointwise residual evaluation on a \emph{verification point set}, distinct from (and typically much larger/denser than) the training collocation set, rather than from a finite-element-style mesh.

\subsubsection{Verification-Point Residual Indicator}
\label{sec:indicator}

Let $\widehat{Q} = \{(x_k,t_k)\}_{k=1}^{K}\subset Q$ and $\widehat{\Sigma} = \{(x_k^b,t_k^b)\}_{k=1}^{K_b}\subset \Sigma$ be dense verification point sets (e.g.\ a fixed low-discrepancy or grid sample, held out from training) with normalized quadrature weights $\{\omega_k\}$ and $\{\omega_k^b\}$, each set summing to $1$. Define the global a posteriori indicator
\begin{equation}
\label{eq:estimator_global}
\eta^2 \;:=\; \sum_{k=1}^{K} \omega_k\,|\mathcal{R}_y(x_k,t_k)|^2
\;+\; \sum_{k=1}^{K} \omega_k\,|\mathcal{R}_\lambda(x_k,t_k)|^2
\;+\; \sum_{k=1}^{K_b} \omega_k^b\,|\mathcal{R}_u(x_k^b,t_k^b)|^2
\;+\; \eta_{\text{ic}}^2,
\end{equation}
where $\eta_{\text{ic}}^2$ approximates $\|Y_\theta(\cdot,0)-y_0\|_{L^2(\Omega)}^2$. The adjoint terminal mismatch is identically zero by \eqref{eq:hard_adjoint}. Thus $\eta$ is evaluated on points independent of training and approximates the continuous residual quantity entering Theorem~\ref{thm:apriori}. A post-hoc partition may be used only to visualize local contributions; it introduces no jump terms.

\subsubsection{A Posteriori Bound}
\label{sec:bound-theorem}

\begin{theorem}[Computable a posteriori bound]
\label{thm:aposteriori}
Under Assumptions (A1)--(A3), and with $\eta$ as in \eqref{eq:estimator_global} computed on a sufficiently fine verification set,
\begin{equation}
\| (e_y, e_\lambda, e_u) \|_{\mathcal{E}} \;\lesssim\; \epsilon_{\rm net}+C_{\rm stab}\bigl(\eta + \mathcal{E}_{\text{quad}}\bigr),
\end{equation}
where $\mathcal{E}_{\text{quad}}$ is the (typically negligible, if $K,K_b$ are large) quadrature error of \eqref{eq:estimator_global} itself relative to the exact $L^2$ residual norms.
\end{theorem}

\begin{proof}[Sketch of proof]
This follows conditionally from Theorem~\ref{thm:apriori} by identifying $\eta$ as a quadrature approximation of the continuous residual norm and retaining both the best-approximation and verification-quadrature contributions.
\end{proof}

\begin{remark}[Contrast with mesh-based estimators]
\label{rem:fem-contrast}
A finite-element a posteriori estimator is typically accompanied by a matching \emph{lower} bound, $\eta(E) \lesssim \|e\|_{1,D_E}$ up to data oscillation (the ``efficiency'' half of the classical reliability--efficiency pair), proved with element-wise bubble functions supported on individual mesh cells. This construction relies essentially on having a mesh with compactly supported, element-local test functions. In the present meshfree PINN setting there is no such local test-function machinery available in general, and we are not aware of a proof of a matching lower bound for collocation-based PINN residual indicators of this type. This is a difference in what the two methods' geometric structure makes available, not a deficiency of \eqref{eq:estimator_global}: $\eta$ should be read as a \emph{one-sided, computable bound} -- useful for monitoring convergence and flagging poorly-resolved regions of $Q$ -- rather than as a certified two-sided (mesh-style) estimator. A two-sided meshfree indicator is left as an open direction.
\end{remark}

\subsection{Practical Computation and Empirical Validation}
\label{sec:practical}

Combining Theorem~\ref{thm:apriori} with the identification of $\eta$ as a quadrature estimate of $\epsilon_{\text{loss}}$, we have informally
\begin{equation}
\| (e_y, e_\lambda, e_u) \|_{\mathcal{E}}^2 \;\lesssim\; \mathcal{L}_{\text{ind}} + \mathcal{E}_{\text{approx}} + \mathcal{E}_{\text{quad}},
\end{equation}
so a small training loss is \emph{necessary but not sufficient} for a small true error: it must be accompanied by (i) a sufficiently expressive architecture ($\mathcal{E}_{\text{approx}}$ small) and (ii) a sufficiently dense/well-distributed collocation set ($\mathcal{E}_{\text{quad}}$ small). In particular, a low training loss evaluated \emph{on the training points themselves} can be misleading if the network overfits those points; this is precisely why $\eta$ is evaluated on an independent verification set (Section~\ref{sec:indicator}).

\subsubsection{Algorithm}

The a posteriori indicator is computed after training as follows:
\begin{enumerate}
    \item Sample a verification point set $\widehat{Q}\cup\widehat{\Sigma}$, independent of the training collocation points (e.g.\ a fresh random draw, or a fixed low-discrepancy grid), together with quadrature weights.
    \item Evaluate $Y_\theta, U_\psi, \Lambda_\phi$ and their required derivatives at these points via automatic differentiation, and compute the residuals $\mathcal{R}_y, \mathcal{R}_\lambda, \mathcal{R}_u$.
    \item Form the weighted sum \eqref{eq:estimator_global}.
    \item (Optional, diagnostic only) Bin the pointwise squared residuals over a post-hoc space-time partition purely for visualization of where the error is concentrated; this partition does not enter the estimator's value.
\end{enumerate}
The computational cost is comparable to one evaluation of the loss function on $K+K_b$ points, i.e.\ negligible relative to training.

\subsubsection{Empirical validation on Examples 1 and 2}
\label{sec:empirical-validation}

Because both numerical tests of Section~4 are manufactured-solution problems, the true error $\|e\|_{\mathcal{E}}$ against $(y^*,u^*,\lambda^*)$ is directly computable, and no separate reference FEM solve is required to validate $\eta$: unlike the general setting of \eqref{eq:decomp}, here we can compare $\eta$ against the true error on the same trained models used in Section~4. We therefore validate the indicator directly on Examples~1 and~2 as follows:
\begin{enumerate}
    \item For each example and each collocation count $N=N_{\rm int}=N_{\rm bc}$, evaluate $\eta$ from \eqref{eq:estimator_global} on a verification set independent of the $N$ training points (the same verification grid used for the residual maps of Section~4).
    \item Compute the true error $\|e\|_{\mathcal{E}}$ against the manufactured solution.
    \item Report the empirical ratio $\eta/\|e\|_{\mathcal{E}}$ (the \emph{effectivity index}) as $N$ varies, rather than assuming a fixed algebraic rate; unlike FEM, PINN convergence in $N$ is not generally guaranteed to follow a clean power law, and any observed rate should be reported empirically.
\end{enumerate}

\begin{table}[H]
\centering
\caption{Effectivity index $\eta/\|e\|_{\mathcal{E}}$ on Examples~1 and~2. \textit{Values to be filled in from the trained checkpoints of Section~4; not yet computed in this draft.}}
\label{tab:effectivity}
\begin{tabular}{llcccc}
\toprule
Example & $N$ & $\|e\|_{\mathcal{E}}$ & $\eta$ & $\eta/\|e\|_{\mathcal{E}}$ \\
\midrule
1 & -- & [TBD] & [TBD] & [TBD] \\
1 & -- & [TBD] & [TBD] & [TBD] \\
2 & -- & [TBD] & [TBD] & [TBD] \\
2 & -- & [TBD] & [TBD] & [TBD] \\
\bottomrule
\end{tabular}
\end{table}

This validates Theorem~\ref{thm:aposteriori} empirically on the present problems; it does not, and is not intended to, validate a lower (efficiency) bound (Remark~\ref{rem:fem-contrast}). For a genuinely low-regularity or non-manufactured problem, where $\|e\|_{\mathcal{E}}$ is unavailable, the same procedure without step 2 is exactly the practical use case for which $\eta$ is intended.

\subsection{Summary of the Error Estimation Framework}
\label{sec:error-summary}

The proposed error estimation framework provides:
\begin{itemize}
    \item \textbf{A conditional a priori bound} (Theorem~\ref{thm:apriori}), valid under explicit well-posedness, regularity, and approximation-capacity assumptions (A1)--(A4).
    \item \textbf{A computable a posteriori bound} \eqref{eq:estimator_global}, built from pointwise residual evaluation on a verification set, with no artificial mesh-jump terms, together with an honest statement that a matching lower (efficiency) bound is not established here for the meshfree setting (Remark~\ref{rem:fem-contrast}).
    \item \textbf{An empirical validation protocol} (Section~\ref{sec:empirical-validation}), applied directly to Examples~1 and~2, reporting the effectivity index $\eta/\|e\|_{\mathcal{E}}$.
\end{itemize}

\begin{remark}
The Dirichlet boundary control setting analyzed here (Assumption (A2)) is a known source of low-regularity difficulties for the classical PDE theory; results here should not be read as removing that difficulty, only as working under regularity hypotheses strong enough to bypass it. Extending the framework to genuinely low-regularity Dirichlet data, and to a two-sided meshfree a posteriori bound, are left for future work.
\end{remark}